
% Here is the abstract.
\begin{abstract}
Il Machine Learning rappresenta uno dei settori più rilevanti dell’informatica moderna, grazie alla sua capacità di estrarre conoscenza dai dati e di supportare processi decisionali in numerosi ambiti applicativi. Negli ultimi anni, il miglioramento delle prestazioni dei modelli di apprendimento automatico è stato ottenuto principalmente attraverso l’aumento della complessità architetturale e della quantità di dati utilizzati in fase di addestramento. Sebbene tale approccio abbia prodotto risultati significativi, esso comporta un incremento dei costi computazionali e una maggiore difficoltà nella gestione dei dataset.

La crescente complessità dei modelli pone problematiche rilevanti in termini di efficienza, scalabilità e impiego delle risorse hardware, soprattutto in contesti in cui tali risorse risultano limitate. Per questo motivo, la ricerca si è orientata verso lo sviluppo di tecniche di ottimizzazione finalizzate a ridurre il carico computazionale dei modelli, mantenendo al contempo un livello di accuratezza accettabile.

Tra le tecniche più diffuse in questo ambito rientra la Knowledge Distillation, un approccio che consente di trasferire la conoscenza da un modello complesso, denominato teacher, a un modello più semplice, denominato student. Tale metodologia permette di ottenere modelli meno onerosi dal punto di vista computazionale, risultando particolarmente vantaggiosa nella fase di inferenza. Tuttavia, il processo di distillazione richiede spesso dataset di grandi dimensioni e una fase di addestramento articolata.

Un approccio alternativo è rappresentato dal Reflection Tuning, una tecnica più recente che si basa su meccanismi di auto-valutazione e miglioramento iterativo delle risposte del modello. A differenza della Knowledge Distillation, il Reflection Tuning non richiede necessariamente la presenza di un modello di riferimento esterno, offrendo potenziali benefici in termini di semplicità del processo di addestramento e di riduzione dei requisiti computazionali.

L’obiettivo di questa tesi è analizzare e confrontare la Knowledge Distillation e il Reflection Tuning, con particolare attenzione alle differenze in termini di complessità computazionale e di dimensione dei dataset necessari per l’addestramento. Attraverso un’analisi comparativa, il lavoro intende evidenziare vantaggi e limiti di ciascun approccio, fornendo un contributo utile alla comprensione delle principali strategie di ottimizzazione dei modelli di Machine Learning

\end{abstract}
%----